\begin{table}[!h]
\centering
\caption{Occupational Transitions: 1940 Elevator Operators in 1950}
\label{tab:transition_matrix}
\centering
\begin{threeparttable}
\begin{tabular}[t]{lll}
\toprule
1950 Occupation & N & Share (\%)\\
\midrule
Not in labor force & 6,330 & 16.4\\
Elevator operator & 6,106 & 15.8\\
Other & 5,191 & 13.5\\
Operative/manufacturing & 4,888 & 12.7\\
Craftsman & 4,690 & 12.2\\
\addlinespace
Farm worker & 4,226 & 11.0\\
Other service worker & 2,015 & 5.2\\
Clerical/sales & 1,894 & 4.9\\
Other building service & 1,801 & 4.7\\
Professional/technical & 1,054 & 2.7\\
\addlinespace
Manager/proprietor & 281 & 0.7\\
Laborer & 86 & 0.2\\
\bottomrule
\end{tabular}
\begin{tablenotes}
\item \textit{Note: }
\item $N = 38{,}562$ linked elevator operators. ``Not in labor force'' uses the census-coded labor force status variable ($N = 6{,}330$); Figure~\ref{fig:occupation_flow} uses a broader NLF definition that includes 141 additional individuals with military or unclassifiable occupation codes ($N = 6{,}471$). Source: IPUMS + MLP v2.0 linked panel.
\end{tablenotes}
\end{threeparttable}
\end{table}
